\subsection{Дизайн експерименту}

У даному підрозділі описано методологію проведення експериментального
дослідження якості машинного перекладу, включаючи процедуру оцінювання,
статистичні методи аналізу та забезпечення валідності результатів.

\subsubsection{Загальна схема експерименту}

Експеримент побудовано за схемою порівняння машинного перекладу з
референсним людським перекладом.
Основні етапи:

\begin{enumerate}
    \item \textbf{Підготовка корпусу:} сегментація аудіозаписів на фрагменти
    оптимальної тривалості (3--5 хвилин);
    \item \textbf{Автоматична обробка:} послідовне застосування ASR (Whisper)
    та MT (OPUS-MT) до кожного сегмента;
    \item \textbf{Створення референсу:} ручна транскрипція та підготовка
    референсних перекладів тих самих сегментів (деталі в розділі~\ref{sec:corpus});
    \item \textbf{Автоматичне оцінювання:} обчислення метрик BLEU, BERTScore,
    METEOR для кожного сегмента;
    \item \textbf{Експертне оцінювання:} оцінка адекватності та плавності,
    анотація помилок;
    \item \textbf{Статистичний аналіз:} агрегація результатів, перевірка
    гіпотез, аналіз кореляцій.
\end{enumerate}

\subsubsection{Процедура автоматичного оцінювання}

Автоматичні метрики обчислюються з використанням стандартизованих
інструментів для забезпечення відтворюваності результатів:

\begin{itemize}
    \item \textbf{BLEU:} бібліотека SacreBLEU~\cite{post_bleu} з токенізацією
    ``13a'' (стандарт WMT);
    \item \textbf{BERTScore:} модель \texttt{microsoft/deberta-xlarge-mnli}
    для англійської мови з налаштуванням \texttt{rescale\_with\_baseline=True};
    \item \textbf{METEOR:} реалізація з пакету NLTK з використанням WordNet
    для синонімів.
\end{itemize}

Оцінювання проводиться на рівні окремих сегментів, що дозволяє:
\begin{itemize}
    \item аналізувати варіативність якості перекладу;
    \item виявляти проблемні ділянки тексту;
    \item досліджувати вплив характеристик вхідного аудіо на якість.
\end{itemize}

\subsubsection{Процедура експертного оцінювання}

Експертне оцінювання зосереджено на двох ключових аспектах: анотації
помилок за типами та оцінці міжфразової когерентності. Оцінювач має
доступ до оригінальної транскрипції українською, машинного перекладу
англійською та референсного перекладу.

\paragraph{Анотація помилок.}
Для детального аналізу проводиться анотація помилок за класифікацією,
наведеною в підрозділі~\ref{subsec:error_class}. Для кожної помилки фіксується:
\begin{itemize}
    \item тип помилки (пропуск, додавання, спотворення тощо);
    \item ступінь серйозності (critical, major, minor);
    \item контекст (фрагмент тексту з помилкою).
\end{itemize}

\subsubsection{Статистичні методи аналізу}

Для аналізу отриманих результатів застосовано наступні статистичні методи:

\paragraph{Описова статистика.}
Для кожної метрики обчислюються:
\begin{itemize}
    \item середнє значення ($\bar{x}$);
    \item стандартне відхилення ($\sigma$);
    \item медіана (Me);
    \item мінімальне та максимальне значення.
\end{itemize}

\paragraph{Кореляційний аналіз.}
Для дослідження зв'язку між автоматичними метриками та людською оцінкою
обчислюється коефіцієнт кореляції Спірмена~\cite{spearman} ($\rho$),
який є більш робастним для ординальних даних:
\begin{equation}
\rho = 1 - \frac{6 \sum d_i^2}{n(n^2 - 1)}
\end{equation}
де $d_i$ — різниця між рангами $i$-го спостереження, $n$ — кількість
спостережень.

\paragraph{Перевірка статистичної значущості.}
Для перевірки значущості відмінностей між групами застосовується критерій
Манна-Уітні (U-тест) як непараметрична альтернатива t-тесту, що не вимагає
припущення про нормальний розподіл:
\begin{equation}
U = n_1 n_2 + \frac{n_1(n_1 + 1)}{2} - R_1
\end{equation}
де $n_1$, $n_2$ — розміри вибірок, $R_1$ — сума рангів першої вибірки.

Рівень значущості встановлено на $\alpha = 0.05$.

\subsubsection{Аналіз впливу помилок ASR на якість перекладу}

Окремим напрямком аналізу є дослідження каскадного ефекту — впливу помилок
автоматичного розпізнавання мовлення (ASR) на якість машинного перекладу.
Для цього:

\begin{enumerate}
    \item Обчислюється WER (Word Error Rate) між ASR-транскрипцією та
    референсною транскрипцією за формулою, наведеною в підрозділі~\ref{subsec:asr}.

    \item Проводиться кореляційний аналіз між WER/CER та метриками якості
    перекладу (BLEU, BERTScore, METEOR).

    \item Виділяються типові помилки ASR, що призводять до критичних
    помилок перекладу (зокрема, спотворення імен власних та абревіатур).
\end{enumerate}

\subsubsection{Забезпечення валідності дослідження}

Для забезпечення валідності та надійності результатів дослідження вжито
наступних заходів:

\begin{itemize}
    \item \textbf{Внутрішня валідність:} використання стандартизованих
    інструментів обчислення метрик, чітких критеріїв експертної оцінки,
    контроль за процедурою оцінювання.

    \item \textbf{Зовнішня валідність:} корпус містить реальні публічні
    промови різної тематики, що підвищує екологічну валідність результатів.

    \item \textbf{Відтворюваність:} всі параметри моделей, версії бібліотек
    та налаштування задокументовані, вихідні дані збережені у структурованому
    форматі.
\end{itemize}

\begin{table}[h]
\centering
\caption{Зведена таблиця методів оцінювання}
\label{tab:evaluation_methods}
\begin{tabular}{|l|l|l|}
\hline
\textbf{Аспект оцінки} & \textbf{Метод} & \textbf{Рівень} \\
\hline
Лексичний збіг & BLEU & Сегмент/корпус \\
Семантична подібність & BERTScore & Сегмент \\
Морфологічний збіг & METEOR & Сегмент \\
Типи помилок & Анотація помилок (MQM) & Сегмент \\
Когерентність & Експертна оцінка & Пара сегментів \\
Якість ASR & WER, CER & Сегмент \\
\hline
\end{tabular}
\end{table}
